\section{The curve, the field and the golden model}
\label{sec:curve}

ECC2K-130 is the Certicom challenge~\cite{certicom1997}
\[
  E:\ y^2 + xy = x^3 + 1 \qquad\text{over }\FF_{2^{131}},
\]
a Koblitz curve~\cite{koblitz1992} defined over $\FF_2$, with group
order $4\cdot\ell$,
\[
  \ell = 680564733841876926932320129493409985129,
\]
a 129-bit prime.\art{ecc2k130/README.md, ``The problems''} The
challenge asks for $k$ with $Q=[k]P$ for published $P$ and $Q$ of
order $\ell$. It has been open since 1997; the 2009 consortium
estimate of the work is $2^{60.9}$ iterations of a
Frobenius-equivariant walk~\cite{bailey2009}, and
Section~\ref{sec:walk-constant} measures this paper's walk at
$2^{60.91}$--$2^{60.92}$ on completed trails.

The curve is identified rather than trusted. $E(\FF_2)$ has four
points, so the Frobenius trace is $t=-1$ and the Koblitz recursion
$s_m = t\,s_{m-1} - 2 s_{m-2}$ gives $\#E(\FF_{2^{131}}) = 2^{131}+1-s_{131}
= 4\ell$. The sibling library's test suite checks the two facts that
make this \emph{this} curve, on random points multiplied by the
cofactor: $\sigma^2+\sigma+2=0$, where $\sigma(x,y)=(x^2,y^2)$ is the
Frobenius, and $\ell\cdot
R=\mathcal{O}$.\art{cryptanalysis:fpga/README.md}\art{cryptanalysis:fpga/model/ecc2k130_test.c}
The published challenge point converts out of the normal basis back
to Certicom's polynomial-basis hex, \texttt{05 1C99BFA6 F18DE467
C80C23B9 8C7994AA}.\art{ecc2k130/README.md, ``Validation''}

\subsection{A type-II optimal normal basis}
\label{sec:onb}

$2\cdot131+1=263$ is prime and $2$ has order $131$ modulo $263$, so
$\FF_{2^{131}}$ has a type-II optimal normal basis. With $\zeta$ a
primitive $263$rd root of unity the basis is
$\gamma_i = \zeta^i + \zeta^{-i}$, $i=1,\ldots,131$, and an element is a
131-bit vector. Every design in this paper, on every platform, is
built on three consequences.\art{ecc2k130/README.md, ``Field representation''}

\begin{center}
\small
\begin{tabular}{@{}p{0.42\textwidth}p{0.48\textwidth}@{}}
\toprule
fact & consequence \\
\midrule
$\gamma_i\gamma_j = \gamma_{i+j} + \gamma_{i-j}$
  (indices folded by $\gamma_{263-k}=\gamma_k$, $\gamma_0=0$)
  & a product is a cyclic convolution of two symmetric 263-bit
    vectors, or, after a Bernstein--Lange change of basis, a
    polynomial product with no reduction step \\
$\gamma_i^2 = \gamma_{2i \bmod 263}$
  & squaring is the index permutation $i\mapsto\mathrm{fold}(2i)$:
    free when loops are unrolled, wiring on an FPGA, so every
    Frobenius power $\sigma^k$ costs nothing \\
the Hamming weight of the vector is invariant under $\sigma$ and under
  negation $-(x,y)=(x,x+y)$
  & the iteration's branch selector and the distinguished-point
    predicate are functions on $\langle\sigma,-1\rangle$-classes that
    read the vector the hardware already holds \\
\bottomrule
\end{tabular}
\end{center}

\subsection{The golden model}
\label{sec:golden}

None of this is taken on trust. The code generator models the field
as the symmetric vectors over $\ZZ/263\ZZ$ modulo the all-ones vector,
with multiplication as cyclic convolution: a complete and obviously
correct model of $\FF_{2^{131}}$, and the one every generated routine is checked against on 64 random
inputs before it is written out.\art{ecc2k130/codegen/field.py}\art{ecc2k130/codegen/gen.py} The bitsliced multiply, square,
inverse, $\sigma^7$ and Hamming weight are differentially tested
against a reference implementation that shares no code with them;
the iteration function is checked to commute with Frobenius and with
negation; the collision solver recovers a planted logarithm; the
distinguished-point comparison is checked exhaustively over all
weights and cutoffs. \texttt{make test} runs 52 such checks across
$\mathrm{GF}(2^{23})$, $\mathrm{GF}(2^{41})$, $\mathrm{GF}(2^{83})$ and
$\mathrm{GF}(2^{131})$.\art{ecc2k130/README.md, ``Validation''}

The same model is the oracle for every other platform. The VHDL
engine's reference script imports it directly and proves the
basis-change maps, the multiplier route, $\sigma^j$, the inversion
chain and a full step before emitting test vectors
(Section~\ref{sec:fpga}).\art{hdl/ecc2k130/ecc2k_ref.py} The sibling
library's independent C model of the same field and curve
(\texttt{cryptanalysis:fpga/model/ecc2k130.c}) was built separately
and re-walks every GPU report that library imports; agreement between
the two models on the same records is equality of 131-bit vectors in
the same permuted basis, not equivalence up to a basis
change~\cite{ecdlppaper}.

Longer runs check the reporting path rather than the algebra. On
$\mathrm{GF}(2^{83})$, $12.6$~M iterations produced 172 distinguished
points, within $30\%$ of the predicted rate, and 60 of them were
recomputed from their seeds by the reference with no mismatch,
including walks that had been restarted. On the challenge curve
itself, with the cutoff loosened from 34 to 50 so that points appear,
$3.1$~M iterations produced $12{,}368$ reports, 25 of which were
recomputed and matched exactly.\art{ecc2k130/README.md, ``Validation''}

\subsection{ECC2K-95 as the end-to-end control}
\label{sec:ecc2k95}

ECC2K-95 is the retired Certicom challenge on the same curve equation
over $\FF_{2^{97}} = \FF_2[t]/(t^{97}+t^6+1)$, group order $4\ell$ with
$\ell = 39614081257132074233778707191$ a 95-bit prime, solved by
Harley and collaborators in 1998~\cite{harley1998}. $2\cdot97+1=195$ is
composite, so there is no type-II optimal normal basis; the client
runs a polynomial-basis backend with the orbit-invariant weight taken
through one generated linear map into a normal basis, which is the
arrangement Harley's own client used.\art{ecc2k130/README.md, ``ECC2K-130 and ECC2K-95''}

The parameters were recovered from Harley's surviving client source
and the generator refuses to emit the curve unless both points lie on
it, both have order $\ell$, and the published answer
\[
  k = 37837308472231540269443981458
\]
satisfies $[k]P=Q$. Expected work is about $2^{44}$ iterations;
Harley's group measured $2.16\cdot10^{13}$ against the
$1.79\cdot10^{13}$ this walk predicts from
$\sqrt{\pi\ell/(4m)}$.\art{ecc2k130/README.md, ``The problems''} Per
iteration the polynomial-basis field is about $1.9\times$ cheaper
than $\mathrm{GF}(2^{131})$ in the normal basis (about $32{,}000$
against $60{,}000$ bitsliced instructions), so ECC2K-95 is a
few GPU-hours where ECC2K-130 is GPU-centuries: a real end-to-end
test of the whole pipeline at a scale that completes.
\art{ecc2k130/README.md, ``ECC2K-95''} The repository's separate
ECC2K-95 CUDA kernel has run only through CPU emulation, where it
recovers 21-bit and 40-bit toy logarithms at $0.94\times$ and
$1.32\times$ its own prediction, and is not a throughput
claim.\art{gpu/ecc2k/README.md}
